\subsubsection{Detekce stěny a určení jejího směru}

V \texttt{Update()} se detekuje, zda se hráč dotýká stěny. Pokud ano a hráč není na zemi, nastaví se \texttt{lastOnWallTime} na \texttt{wallCoyoteTime}. Pokud ne, timer se snižuje o \texttt{Time.deltaTime}. Toto zajišťuje, že hráč má krátký čas na provedení wall jump poté, co opustil stěnu.

Metoda \texttt{DetectWallDirection()} určuje, na které straně stěna je. Existují dva přístupy: buď podle vstupu hráče (pokud hráč drží tlačítko směrem ke stěně, stěna je na této straně), nebo podle otočení postavy jako záloha. První přístup je přesnější, protože umožňuje hráči určit stěnu i když se od ní vzdaluje, druhý přístup slouží jako fallback pro situace bez inputu.